\section{Bochner--Riesz consequences}\label{sec:bochner-riesz}

We apply the established local-smoothing implications of
Beltran--Hickman--Sogge to Corollary~\ref{cor:full-range}. The maximal
result concerns the family in \eqref{eq:intro-br-means}; the subsequent
nonmaximal result concerns the individual linear means.
These means recover a function by gradually admitting higher
frequencies, with order $\delta$ smoothing the boundary of the Fourier
ball. Uniform $L^p$ bounds make each approximation stable. A maximal
bound is stronger: it controls the error for all radii at once, which
allows the elementary convergence for Schwartz data to pass to every
$L^p$ datum by density.

For $p\ge3$, the local-smoothing input has an arbitrarily small
positive loss beyond $1-3/p$. The strict inequality
$\delta>1-3/p$ absorbs that loss in the established implication.
Thus the conclusion below requires only $L^p$ integrability of the
data, in contrast to the Sobolev hypothesis in the maximal half-wave
corollary.

\begin{theorem}[Maximal Bochner--Riesz bounds and almost-everywhere summation]
\label{thm:br-maximal}
Let $3\le p<\infty$ and $\delta>1-3/p$. Then
\begin{equation}\label{eq:br-maximal}
 \norm{B_*^\delta f}_{L^p(\R^3)}
 \le C_{p,\delta}\norm f_{L^p(\R^3)}
 \qquad(f\in\mathcal S(\R^3)).
\end{equation}
The family has a bounded maximal extension to $L^p(\R^3)$, and in that
extension every $f\in L^p(\R^3)$ satisfies
\begin{equation}\label{eq:br-ae}
 \lim_{R\to\infty}B_R^\delta f(x)=f(x)
 \quad\text{for almost every }x\in\R^3.
\end{equation}
In particular, every $\delta>0$ is allowed when $p=3$.
\end{theorem}

\begin{proof}
For $p\ge3$, the sharp fixed-time Sobolev loss is $1-2/p$.
Corollary~\ref{cor:full-range} permits every exponent
$\alpha=1-3/p+\eps$, $\eps>0$, and hence gives the $1/p-$ local
smoothing hypothesis for the Euclidean half-wave in
\cite[Proposition~3.3]{BeltranHickmanSogge2019Survey}. That proposition
applies for $2n/(n-1)\le p<\infty$ and order above
\[
 \max\left\{n\left|\frac1p-\frac12\right|-\frac12,0\right\}.
\]
For $n=3$ and $p\ge3$ this threshold is $1-3/p$. The arbitrary
positive loss in the local-smoothing input is absorbed by the strict
inequality on $\delta$.

The cited proposition uses the multiplier $(1-t|\xi|)_+^\delta$.
Write $\mathcal L_R^\delta$ for the same means with $t=R^{-1}$ and
$\mathcal L_*^\delta f=\sup_{R>0}|\mathcal L_R^\delta f|$.
It gives the asserted strong bound for $\mathcal L_*^\delta$.
Here is the transfer to the squared-radius convention in
\eqref{eq:intro-br-means}. Choose $\chi\in C_c^\infty((0,\infty))$
equal to one on a neighborhood of $[1/2,1]$, and put
\[
 a(\xi)=\chi(|\xi|)(1+|\xi|)^\delta,\qquad
 b(\xi)=(1-\chi(|\xi|))(1-|\xi|^2)_+^\delta.
\]
Both $a$ and $b$ are smooth and compactly supported on $\R^3$, and
\[
 \left(1-\frac{|\xi|^2}{R^2}\right)_+^\delta
 =a(\xi/R)\left(1-\frac{|\xi|}{R}\right)_+^\delta+b(\xi/R).
\]
The inverse Fourier transforms of $a$ and $b$ are Schwartz functions.
Their dilates are bounded in absolute value by
$C_N R^3(1+R|x|)^{-N}$ for any fixed $N>3$, so the
Hardy--Littlewood maximal operator $\mathcal M$ gives the
pointwise estimate
\[
 B_*^\delta f
 \le C_\delta\bigl(\mathcal M(\mathcal L_*^\delta f)+\mathcal Mf\bigr).
\]
Its $L^p$ boundedness for $p>1$ proves \eqref{eq:br-maximal}.

For $g\in\mathcal S(\R^3)$, dominated convergence in the Fourier
integral gives $B_R^\delta g\to g$ uniformly as $R\to\infty$.
For $f\in L^p$, choose $g_j\in\mathcal S$ converging to $f$ in $L^p$
with $\sum_j\norm{g_{j+1}-g_j}_p<\infty$. The maximal inequality
places $\sum_j B_*^\delta(g_{j+1}-g_j)$ in $L^p$. Thus
$B_R^\delta g_j$ converges uniformly in $R$ for almost every $x$,
defining the usual maximal extension; for each fixed $R$ it agrees
with the bounded $L^p$ extension of $B_R^\delta$. For $g\in\mathcal S$,
\[
 \limsup_{R\to\infty}|B_R^\delta f-f|
 \le B_*^\delta(f-g)+|f-g|
 \quad\text{almost everywhere}.
\]
Chebyshev's inequality and \eqref{eq:br-maximal} bound the measure
where this limsup exceeds $\eta>0$ by
$C_{p,\delta}\eta^{-p}\norm{f-g}_p^p$. Density then proves
\eqref{eq:br-ae}.
\end{proof}

\begin{corollary}[The full strict nonmaximal range]
\label{cor:br-nonmaximal}
For $1\le p\le\infty$, with $1/\infty=0$, set
\[
 \delta_{\mathrm{BR}}(p)
 =\max\left\{3\left|\frac1p-\frac12\right|-\frac12,0\right\}.
\]
If $\delta>\delta_{\mathrm{BR}}(p)$, the operators $B_R^\delta$
extend boundedly to $L^p(\R^3)$ with a bound uniform in $R>0$:
\begin{equation}\label{eq:br-nonmaximal}
 \norm{B_R^\delta f}_{L^p(\R^3)}
 \le C_{p,\delta}\norm f_{L^p(\R^3)}.
\end{equation}
\end{corollary}

\begin{proof}
Proposition~3.2 of \cite{BeltranHickmanSogge2019Survey}, with the same
local-smoothing input and the same smooth-factor transfer, gives
\eqref{eq:br-nonmaximal} for $3\le p<\infty$ and
$\delta>1-3/p$; dilation makes the constants uniform in $R$.
For $2\le p\le3$, fix any $\delta>0$ and interpolate the $p=3$
bound at that order with the $L^2$ bound from Plancherel. This gives
the required threshold $\delta_{\mathrm{BR}}(p)=0$ in this interval.
The multiplier is real and even, so duality gives the range
$1<p<2$, since $\delta_{\mathrm{BR}}(p)=\delta_{\mathrm{BR}}(p')$.
Finally, at $p=1$ and $p=\infty$ the stated condition is $\delta>1$.
The inverse Fourier transform $K^\delta$ of $(1-|\xi|^2)_+^\delta$ obeys
$|K^\delta(x)|\le C_\delta(1+|x|)^{-2-\delta}$ and is therefore
integrable at these orders. Its dilates have the same $L^1$ norm,
so Young's inequality supplies both endpoint bounds.
\end{proof}

For $1\le p<\infty$ in this strict range, the uniform bounds also give
$B_R^\delta f\to f$ in $L^p$ as $R\to\infty$. One first checks
this for Schwartz functions with compact Fourier support and then
uses their density in $L^p$ together with the uniform operator bound.

The interpolation and duality in this corollary apply to the
fixed-radius linear operators; the maximal conclusion remains the
range of Theorem~\ref{thm:br-maximal}. For comparison, the separate
article \cite[Theorem~1.1 and Corollary~12.4]{OpenAIBochnerRiesz2026}
proves positive-order $L^3$ boundedness and the same strict nonmaximal
range by a different argument. That article is not an input to the
deduction above.

\subsection{Spherical restriction and Kakeya maximal estimates}
\label{subsec:restriction-kakeya}

The strict multiplier range has further established consequences for
Fourier extension and tube geometry. Let $d\sigma$ be surface measure
on $S^2$, and define the spherical extension operator by
\[
 Eg(x)=\int_{S^2}e^{ix\cdot\omega}g(\omega)\,d\sigma(\omega).
\]
Tao's Bochner--Riesz-to-restriction implication
\cite{Tao1999Restriction}, also described in
\cite[Section~3.1]{BeltranHickmanSogge2019Survey}, applies to
Corollary~\ref{cor:br-nonmaximal}. By duality it gives
\begin{equation}\label{eq:sphere-extension}
 \norm{Eg}_{L^r(\R^3)}\le C_r\norm{g}_{L^r(S^2,d\sigma)},
 \qquad 3<r<\infty.
\end{equation}
Thus the multiplier bounds also control the physical-space
concentration of waves with Fourier support on the sphere. The strict
family of positive-order Bochner--Riesz bounds supplies the premise;
no order-zero multiplier assertion is needed.

For $0<h<1$ and $\omega\in S^2$, let $\mathcal T_h(\omega)$ be
the family of all unit-length tubes of radius $h$ parallel to $\omega$,
at arbitrary spatial positions, and put
\[
 K_hF(\omega)=\sup_{T\in\mathcal T_h(\omega)}
              \frac1{|T|}\int_T|F(x)|\,dx.
\]
The restriction-to-Kakeya implication of
\cite[Section~4, Proposition~4.1]{Wolff1999Kakeya} applies to the
diagonal extension range \eqref{eq:sphere-extension}. For $r>3$
sufficiently close to $3$, it first gives restricted weak type $(a,a)$,
where $a=r/(r-2)<3$, with norm $O(h^{-2(3/a-1)})$ for the tubes above.
Interpolation with the trivial $L^\infty$ bound gives strong type
$(3,3)$ with loss $h^{-2(1-a/3)}$. Since $a\uparrow3$ as
$r\downarrow3$, this yields the Kakeya maximal estimate, for every
$F\in L^3(\R^3)$,
\begin{equation}\label{eq:kakeya-maximal}
 \norm{K_hF}_{L^3(S^2,d\sigma)}
 \le C_\eps h^{-\eps}\norm{F}_{L^3(\R^3)},
 \qquad \eps>0,\quad 0<h<1.
\end{equation}
This quantitatively controls averages over thin tubes in every
direction. In particular, a compact set containing a unit line segment
in each direction has Hausdorff dimension three. Wang--Zahl already
proved the Hausdorff and Minkowski dimension assertions
\cite[Theorem~1.1]{WangZahl2025Kakeya}; the maximal estimate is a
stronger quantitative conclusion. The deductions here use the
established Bochner--Riesz-to-restriction-to-Kakeya chain and also apply
to the independent Bochner--Riesz result cited above.
